\documentclass[aspectratio=169]{beamer}

\usetheme{Madrid}
\usecolortheme{default}

\usepackage[T1]{fontenc}
\usepackage[utf8]{inputenc}
\usepackage[french]{babel}
\usepackage{amsmath,amssymb}
\usepackage{tikz}
\usepackage{graphicx}
\usepackage{xcolor}
\usepackage{booktabs}
\usetikzlibrary{arrows.meta, positioning, decorations.pathreplacing, calc}

\definecolor{cprf-navy}{RGB}{30,39,97}
\definecolor{cprf-ice}{RGB}{202,220,252}
\definecolor{cprf-accent}{RGB}{240,84,84}
\definecolor{cprf-light}{RGB}{235,240,255}
\definecolor{cprf-gray}{RGB}{100,110,140}
\definecolor{cprf-green}{RGB}{40,130,80}

\setbeamercolor{palette primary}{bg=cprf-navy, fg=white}
\setbeamercolor{palette secondary}{bg=cprf-navy!80, fg=white}
\setbeamercolor{palette tertiary}{bg=cprf-navy!60, fg=white}
\setbeamercolor{palette quaternary}{bg=cprf-navy, fg=white}
\setbeamercolor{structure}{fg=cprf-navy}
\setbeamercolor{block title}{bg=cprf-navy, fg=white}
\setbeamercolor{block body}{bg=cprf-light}
\setbeamercolor{block title alerted}{bg=cprf-accent, fg=white}
\setbeamercolor{block body alerted}{bg=cprf-accent!10}
\setbeamercolor{block title example}{bg=cprf-green, fg=white}
\setbeamercolor{block body example}{bg=cprf-green!10}
\setbeamercolor{alerted text}{fg=cprf-accent}
\setbeamercolor{frametitle}{bg=cprf-navy, fg=white}
\setbeamercolor{item}{fg=cprf-navy}

\setbeamertemplate{navigation symbols}{}
\setbeamertemplate{itemize item}{\color{cprf-accent}\textbullet}
\setbeamertemplate{itemize subitem}{\color{cprf-gray}$\circ$}

\makeatletter
\setbeamertemplate{footline}{%
  \leavevmode%
  \hbox{%
    \begin{beamercolorbox}[wd=.5\paperwidth,ht=2.5ex,dp=1ex,left,leftskip=1em]{palette primary}%
      \usebeamerfont{author in head/foot}\insertshortauthor%
    \end{beamercolorbox}%
    \begin{beamercolorbox}[wd=.4\paperwidth,ht=2.5ex,dp=1ex,center]{palette secondary}%
      \usebeamerfont{title in head/foot}\insertshorttitle%
    \end{beamercolorbox}%
    \begin{beamercolorbox}[wd=.1\paperwidth,ht=2.5ex,dp=1ex,right,rightskip=1em]{palette tertiary}%
      \insertframenumber/\inserttotalframenumber%
    \end{beamercolorbox}%
  }%
}
\makeatother

\title[CPRF, Revisited]{Constrained Pseudorandom Functions, Revisited}
\subtitle{Soutenance de projet -- Master 1 Cryptologie}
\author[Mouthon \& Razafinony]{Julian Mouthon \quad Mario Razafinony}
\institute{Sorbonne Universite}
\date{\today}

% ============================================================
\begin{document}

% ---- SLIDE TITRE ----
{
\setbeamertemplate{footline}{}
\begin{frame}
  \vspace{0.8cm}
  \begin{center}
    {\color{cprf-navy}\LARGE\textbf{Constrained Pseudorandom Functions}}\\[0.4cm]
    {\color{cprf-gray}\large Revisited}\\[1.0cm]
    \begin{beamercolorbox}[rounded=true,sep=10pt,center,wd=0.75\textwidth]{block title}
      \large Soutenance de projet --- Master 1 Cryptologie
    \end{beamercolorbox}
    \vspace{0.7cm}
    {\large Julian Mouthon \quad \& \quad Mario Razafinony}\\[0.3cm]
    {\color{cprf-gray} Sorbonne Universite --- \today}
  \end{center}
\end{frame}
}

% ---- PLAN ----
\begin{frame}{Plan de la presentation}
  \tableofcontents
\end{frame}

% ============================================================
\section{Motivations et definitions}

\begin{frame}{Pourquoi les CPRF ?}
  \begin{columns}
    \begin{column}{0.55\textwidth}
      \textbf{Probleme concret}
      \begin{itemize}
        \item On veut deleguer l'evaluation d'une PRF sur un sous-ensemble d'entrees
        \item Sans reveler la cle secrete complete
        \item Applications : \textit{Searchable Symmetric Encryption}, stockage externalise
      \end{itemize}
      \vspace{0.5cm}
      \textbf{Solution : les CPRF}
      \begin{itemize}
        \item Une \textbf{PRF} : sortie indiscernable du hasard
        \item Une \textbf{CPRF} : restreint l'evaluation par une contrainte $C$
        \item La cle contrainte $K_C$ ne permet pas de retrouver $K$
      \end{itemize}
    \end{column}
    \begin{column}{0.42\textwidth}
      \begin{block}{Analogie}
        Comme un \textbf{badge d'acces partiel} : il ouvre certaines portes, pas toutes, sans reveler le passe maitre.
      \end{block}
      \vspace{0.4cm}
      \begin{alertblock}{Notre contrainte principale}
        \[ C_n(x) = [x < n] \]
        Autoriser les entrees $< n$.
      \end{alertblock}
    \end{column}
  \end{columns}
\end{frame}

\begin{frame}{Definition formelle d'une CPRF}
  \begin{columns}[t]
    \begin{column}{0.23\textwidth}
      \begin{block}{$\textsf{Setup}(1^\lambda) \to K$}
        Genere la \textbf{cle maitre} $K$
      \end{block}
    \end{column}
    \begin{column}{0.23\textwidth}
      \begin{block}{$\textsf{Constrain}(K, C) \to K_C$}
        Cree une \textbf{cle contrainte} pour $C$
      \end{block}
    \end{column}
    \begin{column}{0.23\textwidth}
      \begin{block}{$\textsf{Eval}(K, x) \to y$}
        Evalue avec la \textbf{cle maitre}
      \end{block}
    \end{column}
    \begin{column}{0.23\textwidth}
      \begin{block}{$\textsf{EvalC}(K_C, x) \to y$}
        Evalue avec $K_C$ si $C(x)=1$
      \end{block}
    \end{column}
  \end{columns}

  \vspace{0.5cm}

  \begin{columns}
    \begin{column}{0.48\textwidth}
      \begin{alertblock}{Propriete de correction}
        \[ \forall x \text{ tel que } C(x)=1 : \]
        \[ \textsf{EvalC}(K_C, x) \;=\; \textsf{Eval}(K, x) \]
        La cle contrainte donne le meme resultat que la cle maitre sur les entrees autorisees.
      \end{alertblock}
    \end{column}
    \begin{column}{0.48\textwidth}
      \begin{exampleblock}{Securite intuitive}
        Connaitre $K_C$ ne doit \textbf{pas} permettre de calculer $\textsf{Eval}(K, x)$ pour $x$ tel que $C(x) = 0$.\\[0.3cm]
        Autrement dit : la cle contrainte ne revele rien sur les entrees interdites.
      \end{exampleblock}
    \end{column}
  \end{columns}
\end{frame}

% ============================================================
\section{Construction BMO17}

\begin{frame}{La construction BMO17 --- Principe}
  \begin{columns}
    \begin{column}{0.48\textwidth}
      \textbf{Idee centrale}
      \begin{itemize}
        \item Permutation a trappe $\pi$ basee sur RSA
        \item Etat initial secret $ST_0$
        \item Evaluer en $c$ = appliquer $\pi^{-1}$ exactement $c$ fois
      \end{itemize}
      \vspace{0.4cm}
      \textbf{Cle maitre :} $(ST_0,\; SK)$
      \[
        \textsf{Eval}\bigl((SK, ST_0),\, c\bigr) = \pi_{SK}^{-c}(ST_0)
      \]
    \end{column}
    \begin{column}{0.50\textwidth}
      \centering
      \begin{tikzpicture}[node distance=1.9cm]
        \tikzset{st/.style={draw=cprf-navy, fill=cprf-light, rectangle, rounded corners,
                            minimum height=8mm, minimum width=14mm, font=\small}}
        \node[st] (s0) {$ST_0$};
        \node[st, right of=s0] (s1) {$ST_1$};
        \node[right of=s1, font=\large] (dots) {$\cdots$};
        \node[st, right of=dots] (sc) {$ST_c$};
        \draw[->, thick, cprf-accent] (s0) -- node[above]{\footnotesize$\pi^{-1}$} (s1);
        \draw[->, thick, cprf-accent] (s1) -- node[above]{\footnotesize$\pi^{-1}$} (dots);
        \draw[->, thick, cprf-accent] (dots) -- node[above]{\footnotesize$\pi^{-1}$} (sc);
        \node[below=0.3cm of s0, font=\footnotesize, color=cprf-gray] {etat initial};
        \node[below=0.3cm of sc, font=\footnotesize, color=cprf-navy, align=center]
              {resultat\\$= F(K,c)$};
      \end{tikzpicture}
    \end{column}
  \end{columns}
\end{frame}

\begin{frame}{La construction BMO17 --- Cle contrainte}
  \begin{columns}
    \begin{column}{0.50\textwidth}
      \textbf{Cle contrainte pour $C(x) = [x < n]$}
      \begin{itemize}
        \item Cle publique $PK$
        \item Etat $ST_n = \pi_{SK}^{-n}(ST_0)$
        \item Seuil $n$
        \item[$\Rightarrow$] La cle secrete $SK$ n'est \textbf{jamais} incluse
      \end{itemize}
      \vspace{0.4cm}
      \textbf{Evaluation contrainte} ($c < n$) :
      \[
        \textsf{EvalC}\bigl((PK, ST_n, n),\, c\bigr) = \pi_{PK}^{\,n-c}(ST_n)
      \]
    \end{column}
    \begin{column}{0.48\textwidth}
      \centering
      \begin{tikzpicture}[node distance=1.9cm]
        \tikzset{st/.style={draw=cprf-navy, fill=cprf-light, rectangle, rounded corners,
                            minimum height=8mm, minimum width=14mm, font=\small}}
        \node[st] (sn) {$ST_n$};
        \node[st, right of=sn] (sm) {$\cdots$};
        \node[st, right of=sm] (sc) {$ST_c$};
        \draw[->, thick, cprf-navy] (sn) -- node[above]{\footnotesize$\pi$} (sm);
        \draw[->, thick, cprf-navy] (sm) -- node[above]{\footnotesize$\pi$} (sc);
        \node[below=0.3cm of sn, font=\footnotesize, color=cprf-gray, align=center]
              {cle\\contrainte};
        \node[below=0.3cm of sc, font=\footnotesize, color=cprf-navy, align=center]
              {resultat\\$= F(K,c)$};
      \end{tikzpicture}
      \vspace{0.4cm}
      \begin{alertblock}{Correction}
        $\pi^{n-c}(ST_n) = \pi^{n-c}(\pi^{-n}(ST_0)) = \pi^{-c}(ST_0)$ \checkmark
      \end{alertblock}
    \end{column}
  \end{columns}
\end{frame}

% ============================================================
\section{L'attaque CJ25}

\begin{frame}{L'attaque CJ25 --- Intuition}
  \begin{columns}
    \begin{column}{0.50\textwidth}
      \textbf{Idee de l'attaque}
      \begin{enumerate}
        \item Demander la cle contrainte $(PK, ST_n, n)$
        \item Interroger l'oracle sur $x > n$, obtenir $ST_x$
        \item Calculer $\pi_{PK}^{x-n}(ST_x)$ et comparer a $ST_n$
        \item Si egal $\Rightarrow$ c'est une CPRF !
      \end{enumerate}
      \vspace{0.4cm}
      \begin{alertblock}{Avantage de l'attaquant}
        \[\Pr[\mathcal{A} \text{ distingue}] \approx 1 - \tfrac{1}{N}\]
        Negligeable seulement si $N$ est exponentiel.
      \end{alertblock}
    \end{column}
    \begin{column}{0.48\textwidth}
      \centering
      \begin{tikzpicture}[node distance=1.7cm]
        \tikzset{st/.style={draw=cprf-navy, fill=cprf-light, rectangle, rounded corners,
                            minimum height=7mm, minimum width=12mm, font=\small}}
        \node[st] (s0) {$ST_0$};
        \node[st, right of=s0] (sn) {$ST_n$};
        \node[st, right of=sn] (sn1) {$ST_{n+1}$};
        \node[st, right of=sn1] (sx) {$ST_x$};
        \draw[->, thick] (s0) -- node[above]{\tiny$\pi$} (sn);
        \draw[->, thick] (sn) -- node[above]{\tiny$\pi$} (sn1);
        \draw[->, thick] (sn1) -- node[above]{\tiny$\pi$} (sx);
        \draw[decorate, decoration={brace,mirror,amplitude=4pt}, cprf-navy, thick]
          (s0.south west) -- (sn.south east)
          node[midway, below=7pt, font=\scriptsize, cprf-navy] {contrainte};
        \draw[decorate, decoration={brace,mirror,amplitude=4pt}, cprf-accent, thick]
          (sn1.south west) -- (sx.south east)
          node[midway, below=7pt, font=\scriptsize, cprf-accent] {hors contrainte};
        \draw[->, thick, dashed, cprf-accent]
          (sx.south) .. controls +(0,-1.4) and +(0,-1.4) ..
          node[below, font=\scriptsize, cprf-accent, align=center]{verif. $\pi^{x-n}$} (sn.south);
      \end{tikzpicture}
    \end{column}
  \end{columns}
\end{frame}

\begin{frame}{L'attaque CJ25 --- Implementation}
  \begin{columns}
    \begin{column}{0.50\textwidth}
      \textbf{Architecture client--serveur TCP}
      \vspace{0.2cm}
      \begin{block}{Serveur = Oracle}
        \begin{itemize}
          \item \texttt{EVAL(x)} : retourne $F(K,x)$ ou valeur aleatoire selon le monde
          \item \texttt{CONSTRAIN(n)} : retourne $(PK, ST_n, n)$
        \end{itemize}
      \end{block}
      \vspace{0.3cm}
      \begin{block}{Client = Attaquant}
        \begin{itemize}
          \item Choisit $n$, demande la cle contrainte
          \item Fait des requetes $x > n$
          \item Verifie la coherence via $\pi^{x-n}$
        \end{itemize}
      \end{block}
    \end{column}
    \begin{column}{0.46\textwidth}
      \begin{alertblock}{Resultat}
        L'attaque reussit avec probabilite $\approx 1$ des la premiere requete sur la version \textbf{non hachee}.
      \end{alertblock}
      \vspace{0.4cm}
      \begin{exampleblock}{Conclusion}
        BMO17 sans hachage \textbf{n'est pas} IND-securisee. La structure algebrique est completement exploitable.
      \end{exampleblock}
    \end{column}
  \end{columns}
\end{frame}

% ============================================================
\section{Alternatives a RSA}

\begin{frame}{Alternatives a la permutation RSA}
  \begin{columns}[t]
    \begin{column}{0.32\textwidth}
      \begin{block}{$F^{\text{WEAK}}$ --- lineaire}
        \[ \text{Eval}(S, \vec{x}) = S \cdot \vec{x} \]
        $S \in (\mathbb{Z}/p\mathbb{Z})^{M \times N}$
      \end{block}
      \vspace{0.15cm}
      \begin{alertblock}{Cassee}
        Systeme lineaire soluble avec $N-1$ requetes : on retrouve $S$ entierement. \textbf{Pas de trappe iterative.}
      \end{alertblock}
    \end{column}
    \begin{column}{0.32\textwidth}
      \begin{block}{Rabin}
        \[ f(x) = x^2 \bmod n \]
        Basee sur la factorisation.
      \end{block}
      \vspace{0.15cm}
      \begin{alertblock}{Non injective}
        4 pre-images par image : impossible d'iterer $f^{-1}$ de facon deterministe. \textbf{Ambiguite a chaque inversion.}
      \end{alertblock}
    \end{column}
    \begin{column}{0.32\textwidth}
      \begin{block}{AES}
        \[ f(x) = E_k(x) \]
        Bijection sur blocs 128 bits.
      \end{block}
      \vspace{0.15cm}
      \begin{alertblock}{Sans trappe publique}
        Inverser necessite $k$ : pas de separation public/prive. \textbf{Pas de cle contrainte possible.}
      \end{alertblock}
    \end{column}
  \end{columns}
  \vspace{0.2cm}
  \begin{block}{}
    \centering
    \textbf{Bijectivite} $+$ \textbf{Trappe public/prive} $+$ \textbf{Efficacite} $\;\Rightarrow\;$ \textbf{Seul RSA satisfait les trois.}
  \end{block}
\end{frame}

\begin{frame}{$F^{\text{WEAK}}$ : attaque et renforcement}
  \begin{columns}
    \begin{column}{0.50\textwidth}
      \textbf{Attaque par systeme lineaire}
      \begin{itemize}
        \item Choisir $\vec{y} = (0,\ldots,0,1)$, obtenir $S_{\vec{y}}$
        \item Calculer $k_i = S_{\vec{y}_i} \cdot \vec{x}_j$ pour $N-1$ vecteurs $\vec{x}_j \perp \vec{y}$
        \item Resoudre $X \cdot \vec{S}_i^T = \vec{k}_i$ : systeme $(N-1)\times(N-1)$ inversible
        \item[$\Rightarrow$] On retrouve $S$ en entier avec 1 requete d'evaluation
      \end{itemize}
      \vspace{0.3cm}
      \begin{alertblock}{$F^{\text{WEAK}}$ est cassee}
        Un seul appel a \texttt{EVAL} suffit a retrouver la cle maitre.
      \end{alertblock}
    \end{column}
    \begin{column}{0.46\textwidth}
      \begin{block}{Renforcement par hachage}
        \[ \text{Eval}_H(\vec{x}) = H(S \cdot \vec{x}) \]
        Retrouver $S$ necessite une pre-image de $H$ $\Rightarrow$ infaisable.
      \end{block}
      \vspace{0.3cm}
      \begin{exampleblock}{Mais une limite demeure}
        $F^{\text{WEAK}}$ n'est \textbf{pas} IND$^*$-NC-CPRF securisee. Le theoreme de CJ25 ne peut donc pas s'appliquer directement pour prouver la securite de $\textsf{Hashed}[F^{\text{WEAK}}]$.
      \end{exampleblock}
    \end{column}
  \end{columns}
\end{frame}

% ============================================================
\section{CPRF hachee et securite}

\begin{frame}{Renforcement par hachage --- Principe}
  \begin{columns}
    \begin{column}{0.50\textwidth}
      \textbf{Idee} : appliquer un oracle aleatoire $P$ sur la sortie de la CPRF
      \[
        \textsf{Eval}_H(K, x) = P\!\bigl(\textsf{Eval}(K, x)\bigr)
      \]
      \vspace{0.3cm}
      \textbf{Lazy Sampling} : $P$ attribue une valeur aleatoire a chaque entree, memorisee
      \begin{itemize}
        \item Si $z \notin \sigma_P$ : choisir $y \leftarrow \{0,1\}^\lambda$, memoriser $(z,y)$
        \item Sinon : retourner le $y$ memorise
      \end{itemize}
      \vspace{0.3cm}
      \textbf{Deux implementations}
      \begin{itemize}
        \item Version SHA-256 (pratique)
        \item Version lazy sampling (modele theorique CJ25)
      \end{itemize}
    \end{column}
    \begin{column}{0.48\textwidth}
      \centering
      \begin{tikzpicture}[node distance=2.0cm]
        \tikzset{
          st/.style={draw=cprf-navy, fill=cprf-light, rectangle, rounded corners,
                     minimum height=8mm, minimum width=16mm, font=\small},
          hst/.style={draw=cprf-accent, fill=cprf-accent!10, rectangle, rounded corners,
                      minimum height=8mm, minimum width=16mm, font=\small}
        }
        \node[st] (a0) {$ST_x$};
        \node[st, right=2.2cm of a0] (a1) {$ST_n$};
        \node[above=0.15cm of a0, font=\small\bfseries, cprf-navy] {CPRF};
        \draw[->, thick, cprf-navy] (a0) -- node[above]{\small$\pi^{x-n}$} (a1);

        \node[hst, below=1.5cm of a0] (b0) {$H(ST_x)$};
        \node[hst, right=2.2cm of b0] (b1) {$H(ST_n)$};
        \node[above=0.15cm of b0, font=\small\bfseries, cprf-accent] {Hashed CPRF};
        \draw[dashed, thick, cprf-accent] (b0) -- node[above]{\Large$\times$} (b1);

        \draw[->, cprf-gray, thick] (a0) -- node[left]{\small$H$} (b0);
        \draw[->, cprf-gray, thick] (a1) -- node[right]{\small$H$} (b1);
      \end{tikzpicture}
      \vspace{0.3cm}
      \begin{alertblock}{Effet cle}
        L'attaquant ne peut plus verifier $\pi^{x-n}(ST_x) = ST_n$ car il ne voit que les haches.
      \end{alertblock}
    \end{column}
  \end{columns}
\end{frame}

\begin{frame}{Renforcement par hachage --- Theoreme CJ25}
  \begin{columns}
    \begin{column}{0.52\textwidth}
      \begin{block}{Theoreme (Cheng--Jaeger 2025)}
        Si les trois conditions suivantes sont satisfaites :
        \begin{enumerate}
          \item $\min_x |\textsf{CF.Out}(\lambda, x)|$ est super-polynomial en $\lambda$
          \item $CF$ est \textbf{IND$^*$-NC-CPRF securisee}
          \item L'oracle aleatoire est \textbf{non-programmable}
        \end{enumerate}
        \vspace{0.2cm}
        Alors $\textsf{Hashed}[CF]$ est \textbf{SIM$^*$-AC-CPRF securisee}.
      \end{block}
    \end{column}
    \begin{column}{0.44\textwidth}
      \textbf{Application a BMO17 (RSA)}
      \begin{itemize}
        \item \checkmark Espace de sortie $\approx 2^{4096}$ : super-polynomial
        \item \checkmark RSA est IND$^*$-NC securisee
        \item \checkmark Oracle non-programmable (lazy sampling)
      \end{itemize}
      \vspace{0.3cm}
      \begin{exampleblock}{Bilan}
        $\textsf{Hashed}[\text{BMO17}]$ est \textbf{prouvablement securisee} contre l'attaque CJ25.
      \end{exampleblock}
      \vspace{0.2cm}
      \textbf{Application a $F^{\text{WEAK}}$}\\
      $\times$ Condition (2) non satisfaite\\
      $\Rightarrow$ Le theoreme ne s'applique pas.
    \end{column}
  \end{columns}
\end{frame}

% ============================================================
\section{Resultats experimentaux}

\begin{frame}{Resultats experimentaux --- Configuration}
  \begin{columns}
    \begin{column}{0.50\textwidth}
      \textbf{Configuration materielle}
      \vspace{0.3cm}

      \begin{tabular}{ll}
        \toprule
        \textbf{Composant} & \textbf{Detail} \\
        \midrule
        Machine      & Dell Inspiron 3501 \\
        CPU          & Intel Core i7-1165G7 \\
                     & 4 coeurs, 8 threads \\
        Frequence    & 2.80 GHz (jusqu'a 4.70) \\
        RAM          & 7.5 GB \\
        OS           & Ubuntu 24.04.4 LTS \\
        Noyau        & Linux 6.8.0 \\
        Architecture & x86\_64 \\
        GPU          & Aucun \\
        \bottomrule
      \end{tabular}
    \end{column}
    \begin{column}{0.46\textwidth}
      \textbf{Protocole experimental}
      \begin{itemize}
        \item 100 requetes d'evaluation par variante
        \item Taux de succes de l'attaque CJ25 mesure en fonction du nombre de requetes
        \item 4 variantes testees :
          \begin{itemize}
            \item CPRF BMO17 non hachee
            \item CPRF BMO17 hachee SHA-256
            \item CPRF BMO17 hachee lazy sampling
            \item $F^{\text{WEAK}}$ non hachee
          \end{itemize}
      \end{itemize}
    \end{column}
  \end{columns}
\end{frame}

\begin{frame}{Resultats experimentaux --- Observations}
  \begin{columns}[t]
    \begin{column}{0.48\textwidth}
      \begin{alertblock}{CPRF non hachee}
        Taux $\approx \mathbf{1}$ des la \textbf{1\`ere requete}.\\
        Structure algebrique directement exploitable.
      \end{alertblock}
      \vspace{0.4cm}
      \begin{alertblock}{$F^{\text{WEAK}}$ non hachee}
        Avantage significatif visible rapidement.\\
        Matrice $S$ retrouvee en quelques requetes.
      \end{alertblock}
    \end{column}
    \begin{column}{0.48\textwidth}
      \begin{exampleblock}{CPRF hachee SHA-256}
        Taux $\approx \mathbf{0.5}$ : decision \textbf{aleatoire}.\\
        L'attaquant ne distingue plus rien.
      \end{exampleblock}
      \vspace{0.4cm}
      \begin{exampleblock}{CPRF hachee lazy sampling}
        Taux $\approx \mathbf{0.5}$ egalement.\\
        Resultats identiques a SHA-256.
      \end{exampleblock}
    \end{column}
  \end{columns}

  \vspace{0.4cm}
  \begin{block}{Conclusion}
    \centering
    Le hachage (SHA-256 ou lazy sampling) \textbf{neutralise completement} l'attaque CJ25.\\
    Les deux approches sont \textbf{equivalentes} en pratique.
  \end{block}
\end{frame}

% ============================================================
\section{Conclusion}

\begin{frame}{Conclusion}
  \begin{columns}
    \begin{column}{0.50\textwidth}
      \textbf{Ce que nous avons realise}
      \begin{itemize}
        \item BMO17 avec RSA 4096 bits (OpenSSL)
        \item Attaque CJ25 en client--serveur TCP
        \item Analyse : $F^{\text{WEAK}}$, Rabin, AES
        \item CPRF hachee : SHA-256 et lazy sampling
        \item Etude de $\textsf{Hashed}[F^{\text{WEAK}}]$
      \end{itemize}
      \vspace{0.25cm}
      \textbf{Resultats principaux}
      \begin{itemize}
        \item Le hachage \textbf{neutralise} l'attaque CJ25
        \item RSA : seule permutation a trappe adaptee
        \item $F^{\text{WEAK}}$ cassee mais instructive
      \end{itemize}
    \end{column}
    \begin{column}{0.46\textwidth}
      \begin{exampleblock}{Perspectives}
        \begin{itemize}
          \item Contraintes plus expressives\\(intervalles, circuits)
          \item Constructions : LWE, couplages
          \item Preuve formelle de $\textsf{Hashed}[F^{\text{WEAK}}]$
        \end{itemize}
      \end{exampleblock}
      \vspace{0.3cm}
      \begin{block}{Code source}
        \centering
        \texttt{\small github.com/julian-mtn/bmo17-cprf}
      \end{block}
    \end{column}
  \end{columns}
  \vspace{0.4cm}
  \begin{center}
    \Large\color{cprf-navy}\textbf{Merci pour votre attention !}
    \quad {\normalsize\color{cprf-gray} Questions ?}
  \end{center}
\end{frame}

\end{document}